\section{引言}

在电动汽车和电网规模储能的快速发展驱动下，钠离子电池（SIBs）已成为锂离子电池技术的有力替代方案 [1,2]。钠资源的天然丰度和低成本，加上钠离子电池与现有锂离子电池制造基础设施的工艺兼容性，使钠离子电池具备规模化部署的潜力。然而，实现这一潜力关键在于开发出合适的负极材料，使其同时兼具高可逆容量、优异的倍率性能和长循环稳定性——这一组合目前仍是钠离子电池发展的核心瓶颈 [3]。

在众多负极候选材料中，碳基材料——尤其是硬碳——是最为成熟的选择，但其比容量相对较低（250--300 mAh g$^{-1}$），且在高电流密度下动力学性能迟缓 [4]。合金型负极通过与钠发生合金化反应，可提供显著更高的理论容量。锡（Sn）拥有最高的理论质量比容量（847 mAh g$^{-1}$）且工作电压最低，但其体积膨胀剧烈（超过 400\%），导致活性物质粉化、固体电解质界面膜（SEI）不稳定、初始库仑效率低以及容量快速衰减 [5]。锑（Sb）具有中等比容量（660 mAh g$^{-1}$），其层状结构有利于初始 Na$^{+}$ 扩散，倍率性能相对较好，但仍遭受约 290\% 的严重体积膨胀，且存在毒性和价格波动问题 [6]。在此背景下，铋（Bi）凭借多重优势脱颖而出：中等偏高的质量比容量（385 mAh g$^{-1}$），极高的体积比容量（约 3800 mAh cm$^{-3}$），适中的工作电位（约 0.5 V vs.\ Na$^{+}$/Na），可有效避免钠枝晶沉积风险，同时 (012) 晶面较大的层间距有利于 Na$^{+}$ 扩散 [7,8]。然而，铋负极的实际应用受到其半金属特性导致的较差的电子导电性，以及在 Bi$\leftrightarrow$Na$_{3}$Bi 相变过程中约 250\% 的体积膨胀的双重制约 [9]。近年来，碳封装策略在缓解上述问题方面取得了显著进展：嵌入机械强度高的碳微球中的铋纳米颗粒实现了超快充电（200 A g$^{-1}$ 下 5.5 秒完成充放电），循环超过 12,000 次 [10]；而包覆花状铋微球的氮掺杂碳框架在 30 A g$^{-1}$ 下实现了 368.2 mAh g$^{-1}$ 的可逆容量，在 10,000 次循环后容量保持率为 82\% [11]。这些结果表明，通过导电骨架进行结构限域是稳定铋负极的强大策略，但这些方案依赖于复杂的碳工程路线。

碳包覆策略被广泛用于增强合金型负极的导电性并缓冲体积变化，但碳与活性材料之间的物理粘附力往往有限，在长时间循环中容易发生剥离 [12]。MXenes——尤其是 Ti$_{3}$C$_{2}$T$_{x}$——为铋基复合材料提供了独特优势的骨架平台：接近 10,000 S cm$^{-1}$ 的类金属导电性使电子快速传输成为可能；亲水性表面端基（---O、---OH、---F）可与活性组分形成强界面结合；力学性能强韧的层状结构本身即能容纳钠化/脱钠过程中可观的体积波动 [13,14]。已有研究探索了将铋与 MXene 纳米片结合用于钠离子电池负极 [15]。最近，一种多孔 Ti$_{3}$C$_{2}$T$_{x}$/铋烯自支撑三明治薄膜被报道用于钠存储，直接证明了构建 Bi-MXene 层状结构的可行性 [16]。此外，一种 MXene@Bi$_{2}$S$_{3}$/Mo$_{7}$S$_{8}$ 三明治异质结构在 5.0 A g$^{-1}$ 下实现了 474.9 mAh g$^{-1}$ 的容量，稳定循环超过 1,400 次，进一步印证了三明治型 MXene 构型在钠离子电池负极中的普遍有效性和巨大潜力 [17]。尽管取得了这些进展，一种结构简单的二元 Bi/MXene 复合材料——无需依赖硫化、碳包覆或多组分异质结工程——同时实现高倍率性能和超长循环稳定性，迄今仍未被实现。

本文报道了一种通过简便水热法合成、并辅以热处理还原的二元 Bi/MXene 复合负极。在此独特的"三明治"结构中，铋纳米颗粒被限域在 Ti$_{3}$C$_{2}$T$_{x}$ MXene 层间。MXene 组分发挥双重功能：构建连续的类金属导电网络以实现快速电子传输，同时机械限制铋在钠化过程中的体积膨胀，从而抑制粉化并保持电极完整性。由此制得的 Bi/MXene 复合材料在 10 A g$^{-1}$ 的高电流密度下比容量达 336.9 mAh g$^{-1}$，为 0.1 A g$^{-1}$ 时容量的 94.5\%，初始库仑效率高达 88.1\%。更为突出的是，该电极在 2 A g$^{-1}$ 下经 5,000 次循环后仍保持约 310 mAh g$^{-1}$ 的比容量，对应于 90.4\% 的容量保持率。这些性能指标共同表明，Bi/MXene 三明治结构可同时实现高倍率性能、高可逆容量和长循环耐久性——有效打破了传统铋基负极长期面临的倍率-容量-稳定性三者难以兼得的权衡困境。本文通过变扫速循环伏安测试、电化学阻抗谱和恒电流间歇滴定技术的系统动力学分析，结合循环后扫描电子显微镜和非原位 X 射线衍射及 X 射线光电子能谱表征，系统阐明了其电化学储能机理。
